% sec-09: Document 09, the premises code.
\docpart{Premises Code}{PDAC LLM Oncology Clinical Trial Site Premises Code}

\section{Scope and Zone Classification}

This document states who and what may be where. The premises are divided into
four zones. The zone scheme fixed here is used unchanged by documents 10, 11
and 12.

\begin{mvtable}
\begin{tabularx}{\textwidth}{>{\raggedright\arraybackslash}p{2.0cm} >{\raggedright\arraybackslash}p{3.4cm} >{\raggedright\arraybackslash}p{3.0cm} >{\raggedright\arraybackslash}X}
\headrow \thead{Zone} & \thead{Rooms} & \thead{Who may enter} & \thead{Credential and record}\\
\midrule
Public & Lobby, waiting, public restrooms, participant education room & Anyone, during hours of operation & None; closed-circuit recording retained 90 days \\
Clinical & Exam rooms, infusion positions, phlebotomy, consultation & Participants with an appointment, escorted visitors, clinical staff & Badge for staff; sign-in for participants; visit record \\
Restricted & Procedure suites, sterile core, soiled utility, pharmacy, specimen laboratory & Named staff on the delegation log; a participant during a scheduled procedure & Badge plus role authorization; access log retained six years \\
Machine & Machine room, network closets, audit trail store & Model governance lead, site safety officer, and named information technology staff & Badge plus second factor; two-person entry for any change; access log retained six years \\
\end{tabularx}
\end{mvtable}
\tabcapl{tab:zones}{The four premises zones. A zone with no access record is a
zone with no control, so every row names both the credential and the record, and
the two restricted-zone retention periods match document 03 §5.}

\section{Access Control and Physical Security}

\begin{enumerate}
\item Badges are issued in four classes matching the four zones. A badge grants
  the lowest zone required for the holder's delegated tasks and no higher.

\item A visitor to the clinical zone is escorted at all times by a badged staff
  member. A visitor may not enter the restricted or machine zones.

\item After-hours access, meaning access outside the hours in document 05 §5, is
  granted only for a procedure already in progress or a participant safety
  event, and generates an alert to the director of clinical operations.

\item Access logs for the restricted and machine zones are retained for six
  years from database lock, the period fixed by document 03 §5, so that a
  question about who was present can be answered for as long as a question about
  what the model produced.

\item Closed-circuit recording covers the public zone, the loading area, and
  every restricted zone door, and does not cover the interior of an exam room, a
  procedure suite, or a restroom.
\end{enumerate}

\section{Robot Operating Envelopes and Interlocks}

This section owns the numbers. Document 02 §5 restates them and may not vary
them.

\begin{mvtable}
\begin{tabularx}{\textwidth}{>{\raggedright\arraybackslash}p{4.6cm} >{\raggedright\arraybackslash}p{2.8cm} >{\raggedright\arraybackslash}X}
\headrow \thead{Requirement} & \thead{Value} & \thead{How it is verified}\\
\midrule
Envelope boundary beyond maximum reach & 1.2 meters & Surveyed and floor-marked at commissioning gate C7 \\
Emergency stop reach from any point inside the envelope & 0.8 meters & Measured with a tape at each survey \\
Emergency stop at each suite door & Required & Visual inspection and actuation test \\
Independent interlock channels & 3, at category 3 & Channel-by-channel fault injection \cite{iec80601277} \\
Motion arrest after a stop command & 250 milliseconds & High-speed timing record \cite{iec80601277} \\
Quasi-static force limit, collaborative operation & 80 newtons & Force gauge at the instrument tip \cite{iso134822014} \\
Restart after an emergency stop & Two-person authorization & Log entry with both names, and a re-verification of the register under document 06 §4 \\
Interlock test interval & Monthly, and after any fault & Test record filed to the robot instance register \\
\end{tabularx}
\end{mvtable}
\tabcapl{tab:envelopes}{Eight envelope and interlock requirements. Each is
verified by a measurement rather than by a manufacturer's assurance, which is
the position the ten-gate assurance work established \cite{h2humanoid}.}

The eight robot types and fourteen instances are distributed as follows: two
robotic surgical platforms, one per suite; two limb-stabilization collaborative
instances; one imaging-guided needle placement instance; two intraoperative
ultrasound imaging assistants; two specimen transport and handling instances;
one pharmacy compounding instance; two rehabilitation gait trainers; and two
sterile processing loading instances.

\section{Clean, Sterile, and Soiled Corridors}

Traffic follows the pressure relationships fixed in document 08 §3 and does not
vary from them.

\begin{enumerate}
\item Sterile supply moves from the sterile core into a procedure suite and
  never in the reverse direction.

\item A used instrument tray moves from a procedure suite into soiled utility
  through the soiled corridor, and never through the sterile core.

\item A specimen transport instance moves between a procedure suite and the
  specimen laboratory on a defined path that does not cross the sterile core.
  The path is marked and is part of the envelope survey at commissioning gate
  C7.

\item Where the geometry of the leasehold forces a single corridor to serve both
  directions, a temporal separation is used instead: soiled movement is confined
  to a scheduled window after the last procedure of the day, and the schedule is
  a controlled record.
\end{enumerate}

\section{Waste Streams and Housekeeping}

\begin{mvtable}
\begin{tabularx}{\textwidth}{>{\raggedright\arraybackslash}p{3.2cm} >{\raggedright\arraybackslash}p{2.7cm} >{\raggedright\arraybackslash}p{2.8cm} >{\raggedright\arraybackslash}X}
\headrow \thead{Stream} & \thead{Container} & \thead{Pickup} & \thead{Record}\\
\midrule
General & Standard, lidded & Twice weekly & None required \\
Recyclable & Standard, labeled & Weekly & None required \\
Regulated medical & Red, rigid, lined & Twice weekly, Tuesday and Friday, 07:00 to 09:00 & Manifest retained three years \\
Cytotoxic and pharmaceutical & Yellow, rigid, sealed & Weekly, by the licensed hauler & Manifest retained three years; USP 800 handling \cite{usp797800} \\
Electronic and media & Locked, tamper-evident & On demand & Destruction certificate; media from the machine zone destroyed rather than resold \\
\end{tabularx}
\end{mvtable}
\tabcapl{tab:waste}{Five waste streams. The pickup windows are identical to
document 05 §5, because a municipal operating standard and a premises procedure
that disagree about a delivery window produce a code violation on an ordinary
Tuesday.}

Housekeeping in the restricted zone is performed by trained staff on a written
schedule, using agents compatible with the robot manufacturers' cleaning
instructions. Terminal cleaning of a procedure suite follows every procedure and
is recorded.

\section{Signage, Wayfinding, and Occupant Conduct}

\begin{enumerate}
\item Every zone boundary carries signage stating the zone, the credential
  required, and the name of the responsible role.

\item Each procedure suite door carries the emergency stop location, the
  envelope boundary notice, and the two-person restart requirement.

\item The wayfinding path from the accessible parking stall to the clinical zone
  entrance is the accessible route described in document 10 §3, and is signed at
  each decision point.

\item No person may enter a marked envelope while a robot instance in that
  envelope is energized, except a member of the procedure team performing a
  delegated task. A participant is positioned before energizing and is never
  moved into an energized envelope.

\item Photography and recording are prohibited in the restricted and machine
  zones without written authorization from the director of clinical operations.
\end{enumerate}
